\begin{center}
    \textgt{\Large{2015～2024年秋季の東北地方太平洋沖 \\
    におけるズワイガニ雌の受精嚢重量の変化}} \\
\end{center}

\begin{center}
    \textgt{\normalsize{永尾次郎・藤原邦浩・冨樫博幸・鈴木勇人・森川英祐・時岡　駿・下光利明（水産機構資源研）}}
\end{center}

\vspace{1\baselineskip}

ズワイガニ雌は，交尾時に雄から受け取った射精物を受精嚢と呼ばれる体内器官に貯蔵し，産卵時に射精物中の精子を用いて卵を受精する（Beninger et al. Elner and Paul?）。性成熟に達した雌は，生涯最後の脱皮（=最終脱皮）を行い，脱皮直後に最初の産卵を行う（引用Watson?）。初回産卵後，雌は脱皮することなく1〜2年に1回程度の間隔で産卵を繰り返す（引用）。交尾は最終脱皮直後の軟甲時と，その後の各産卵前の硬甲時にそれぞれ行われると考えられ，各交尾で雄から受け渡された射精物は受精嚢内に層状に蓄積される（引用Sainte-Marie G.?）。卵の受精には，多くの場合，最も最近受け渡された射精物の精子が使われることがアロザイムならびにDNAマイクロサテライトマーカーを用いた親子判定試験により示唆されている（引用Sevegny, Urbani?）。このため，受精嚢内に蓄えられた射精物の重量は，交尾によって増加する一方，産卵ならびに交尾後の時間経過によって減少すると考えられる。その結果，ズワイガニでは，繁殖可能な雌雄の出会いの減少は，交尾機会を減少を通じて，貯蔵射精物の重さを減らす可能性がある。

東北地方太平洋沖（以後，東北海域）では，東日本大震災以降，ズワイガニに対する漁獲圧が低く抑えらているにも関わらず，資源量が低迷を続けている（引用森川?）。東北海域を対象とした底魚類資源量調査（秋季10-11月）では，2015年から各調査点で漁獲された雌の一部を標本として持ち帰り，受精嚢を摘出して重量を測定している。受精嚢重量は受精嚢本体の重量に貯蔵射精物の重量が加わって変動するため，交尾頻度の低下した状態では重量が減少している可能性がある。

本研究では，東北海域のズワイガニを対象として，交尾頻度の指標として雌の受精嚢重量を
利用することで，2015〜2025年の年間および緯度間の推移を調べた。また，各調査点の雌雄密度および性比との関係を検証した。

\section*{材料と方法}

\subsection*{受精嚢重量}
受精嚢重量データは，調査船若鷹丸（692t，水産研究・教育機構所属）を用いた2015〜2025年の底魚類資源量調査（9〜11月）において，着底トロール網で漁獲した雌（合計986尾）から得た（表1）。

各調査で漁獲した雌は船上で冷凍し，水産研究・教育機構水産資源研究所八戸庁舎に持ち帰った。雌は解凍後にノギス（CD67-S20PM, ミツトヨ）を用いて甲幅を測定し，最終脱皮の有無を腹節の目視観察により判別した（吉田 1941）。さらに外科剪刀（115~mm, アズワン）とピンセット（110~mm, FONTAX）を用いて雌を解剖し，卵巣を挟んで左右1対ある受精嚢を摘出した（Fig 1）。摘出した受精嚢の湿重量は電子天秤（ME4002, METTLER TOLEDO）を用いて，左右それぞれ，精度0.01gで秤量した。

\subsection*{調査年，緯度，雌雄密度および性比}
受精嚢の年間および緯度間の推移を調べるため，調査年ならび各調査点の緯度データを使用した。また，受精嚢重量に対する雌雄密度および性比に関係を検証するため，各年の各調査点における最終脱皮後のズワイガニ雌雄の漁獲尾数および曳網面積に基づいて，調査点別の雌雄密度ならびに性比を計算した。本調査における各種データの収集方法，および調査点別密度の計算方法は，下光ら（2026）を参照下さい。

\subsection*{受精嚢重量の変動要因の検証}
受精嚢重量の変動要因の検証には，一般化線形モデルを用いた。応答変数に受精嚢重量（SPTW），説明変数に，調査年(YR)，緯度(LAT)，ズワイガニ雄密度(MDM)，雌密度（F
DM）および性比（RSEX）を用いた。YRはカテゴリー変数，LAT，MDM，FDMおよびRSEXは連続変数である。

mu = b0 + b1 * YR + b2 * LAT + b3 * MDM + b4 * FDM + b5 * RSEX
SPTW ~ normal(mu, sd)

muは線形予測子，~ normal()は誤差が平均mu, 標準偏差sdの正規分布に従うことを示す。解析には，Rを用い，TMBパケージにより各パラメータを推定した。モデル選択に先立ってvifに基づいて多重共線性を確認して使用する変数を絞り込み，赤池情報基準AICを用いてモデル選択した。









\section*{結果}

\section*{考察}


\section*{引用文献}
\parindent = -1.5zw
\begin{indentation}{1.5zw}{0pt}

Elner and Paul
Beninger et al.


阿部晃治. 1965. 釧路沖におけるヒゴロモエビ（\textit{Pandalopsis coccinata} Urita）の研究 第1報 1. 生殖に関する研究. 北海道立水産試験場報告 4: 13--21 


\end{indentation}

\subsection*{質疑応答}

\parindent = -3zw
\begin{indentation}{3zw}{0pt}

	吉川（水産機構資源研）：モデルや生データで年と水温の関係はみられたか？

\end{indentation}
